\documentclass[11pt]{article}

\usepackage[margin=1in]{geometry}
\usepackage{amsmath,amssymb,amsthm,mathtools}
\usepackage{booktabs,longtable,array}
\usepackage{enumitem}
\usepackage{hyperref}
\usepackage{xcolor}
\usepackage{listings}
\usepackage[T1]{fontenc}
\usepackage{lmodern}

\hypersetup{
  colorlinks=true,
  linkcolor=blue,
  citecolor=blue,
  urlcolor=blue,
  pdftitle={Does the Symbolic LM Calculus Canonically Induce the Modal Measurement Theory?},
  pdfauthor={Research synthesis prepared 18 September 2026}
}

\newtheorem{theorem}{Theorem}[section]
\newtheorem{proposition}[theorem]{Proposition}
\newtheorem{corollary}[theorem]{Corollary}
\newtheorem{lemma}[theorem]{Lemma}
\theoremstyle{definition}
\newtheorem{definition}[theorem]{Definition}
\newtheorem{remark}[theorem]{Remark}

\newcommand{\F}{\mathbb{F}}
\newcommand{\B}{\mathcal{B}}
\newcommand{\Sring}{\mathcal{S}_A}
\newcommand{\Poss}{\mathsf{Poss}}
\newcommand{\Supp}{\operatorname{Supp}_A}
\newcommand{\GL}{\operatorname{GL}}
\newcommand{\Id}{\operatorname{Id}}
\newcommand{\vecop}{\operatorname{vec}}
\newcommand{\LM}{\operatorname{LM}}
\newcommand{\Bool}{\mathbf{2}}

\title{\textbf{Does the Symbolic LM Calculus Canonically Induce the Modal Measurement Theory?}\\
\large Adversarial bridge analysis between formula-valued Logical Matrices and the later chain-ring modal/effect framework}
\author{Research synthesis and independent finite verification}
\date{18 September 2026}

\begin{document}
\maketitle

\begin{abstract}
This report asks whether the formula-valued Logical Matrix (LM) calculus developed in
\emph{Operator-Level Boolean Computation with Correspondence Matrices} can be regarded as the
canonical symbolic source of the later modal measurement theory over the finite chain ring
\(A\cong \F_2[u]/(u^4)\). The answer is mixed but mathematically useful. Paper B already proves
that Boolean valuation commutes with LM pairing. That fact extends naturally to arbitrary matrix,
module, and tensor expressions after an explicit scalar enrichment. However, direct Boolean
valuation cannot produce the nilpotent and unit amplitudes of the chain ring: every Boolean-ring
homomorphism into the local ring \(A\) lands in its idempotents \(\{0,1\}\). The correct positive
bridge is therefore not direct valuation but scalar enrichment
\(\B\otimes_{\F_2} A\), followed by valuation specialization
\(v\otimes \Id_A\). In this enriched calculus, each symbolic amplitude has a canonical Boolean
support predicate whose satisfying valuations are exactly those for which the corresponding concrete
ring-valued modal amplitude is nonzero. This yields a precise symbolic-support theorem connecting
formula semantics to modal possibility.

The report also proves negative boundaries. The later rules that effects are rows of reversible
\(A\)-bases and that a nonzero contraction denotes a possible event are additional measurement
structure, not consequences of the Paper-B LM semantics. The original LM state encoding generates
only coordinate effects under valuation and therefore does not generate the complete 192 reversible
projective two-outcome basis family. A division-free branch update
\(J_i^B=B^{-1}P_iB\) is nevertheless available for every reversible modal basis and is compatible
with projective unit rescaling and valuation. Zero divisors still obstruct a closed unrestricted
process theory: nonzero states need not be tensor-closed, and conditioning can take unimodular
preparations to nonunimodular conditional states.

Independent exhaustive checks reproduce the 16-element chain ring, its 8 units and 2 idempotents,
\(|\GL_2(A)|=24576\), the 192 unordered projective reversible bases, 98,304 branch-support checks,
262,144 symbolic-support fiber checks, and the failure of 144 out of 256 ordinary Boolean-ring
matrix products to remain inside the restricted 16-element Paper-B LM family.
\end{abstract}

\tableofcontents
\newpage

\section{Research question and verdict}

The central question is whether the schematic chain
\[
\text{formula-valued LM}
\longrightarrow
\text{Boolean valuation}
\longrightarrow
\text{CM/coefficient representation}
\longrightarrow
\text{state/resource}
\longrightarrow
\text{effect/measurement}
\longrightarrow
\text{logical or modal outcome}
\]
is a genuine mathematical semantics, rather than a collection of later constructions sharing matrix
notation.

\subsection{Executive verdict}

There is a genuine bridge, but the later modal measurement theory is \emph{not} canonically derivable
from Paper B as it currently stands.

The strongest defensible formulation is:

\begin{quote}
Paper B canonically supplies a symbolic Boolean layer and a valuation-natural contraction calculus.
After independently choosing the chain-ring amplitude algebra \(A\), there is a natural scalar-enriched
symbolic extension in which valuation continues to commute with states, effects, contractions, matrix
operations, and tensor products. In that extension, every ring-valued amplitude has a canonical Boolean
support predicate describing exactly the valuations for which the corresponding modal event is nonzero.
However, the rules that effects are rows of reversible \(A\)-bases and that nonzero means possible are
extra measurement structure. They are not consequences of the original formula-valued LM calculus.
\end{quote}

Thus the clean architecture is
\[
\boxed{
\text{Boolean symbolic LM calculus}
\longrightarrow
\text{\(A\)-enriched symbolic calculus}
\overset{\text{valuation}}{\longrightarrow}
\text{\(A\)-valued states/effects}
\overset{\text{support contract}}{\longrightarrow}
\text{modal outcomes}.}
\]

The missing structure consists of at least three independent ingredients:
\begin{enumerate}[label=(\arabic*)]
\item a chosen cyclic/chain-ring multiplication on the four CM coefficients;
\item a chosen admissible measurement family, such as rows of matrices in \(\GL_n(A)\);
\item the support interpretation \(a\neq 0\Longleftrightarrow\) ``possible.''
\end{enumerate}

\section{Primary sources and methodological boundary}

The analysis treated the manuscript \emph{Operator-Level Boolean Computation with Correspondence
Matrices} (``Paper B'') as background rather than rediscovering its results. The primary Paper-B text
read for this investigation includes Section 8.1--8.2 and its proof appendix. In the parsed copy used
for the review, lines 560--669 define the Boolean algebra of formulas modulo equivalence, the
formula-valued LM, positive and general valuation, logical pairing, Proposition 2, and Theorem 3.
Lines 780--874 contain the appendix proof and finite verification scope.

The later modal source read was \texttt{PURE\_BOOLEAN\_MODAL\_CM\_ADDENDUM.tex}. In the
parsed copy used for the review, lines 45--118 explicitly state that the modal nonzero-support
readout is an ``additional contract,'' introduce the chain ring, define states and separability, and
specify reversible effect bases. Lines 440--539 contain the complete 192-basis support classification,
and lines 620--643 identify the remaining multi-round measurement/update problem.

The supplied adversarial-review archive is included verbatim and unpacked in this research package.
It contains the earlier novelty dossier, theorem candidates, search logs, source-audit files, and a large
set of reproduced computational outputs. The current report adds a separate bridge analysis and
fresh independent checks.

\section{Strict type map}

The investigation keeps the following object types separate.

\begin{longtable}{p{0.18\textwidth}p{0.29\textwidth}p{0.23\textwidth}p{0.23\textwidth}}
\toprule
\textbf{Layer} & \textbf{Object} & \textbf{Map} & \textbf{Status / preserved structure} \\
\midrule
\endhead
Binary truth function & \(C\in M_2(\F_2)\) & \(C\mapsto \LM_{X,Y}(C)\) & Injective for freely assignable \(X,Y\); preserves XOR and entrywise Boolean operations; not ordinary matrix multiplication.\\
\addlinespace
Formula-valued LM & \(\LM_{X,Y}(C)\in M_2(\B)\) & entrywise Boolean valuation \(v\) & Canonical for each valuation; preserves Boolean-ring operations.\\
\addlinespace
Valued CM & \(v(\LM)\in M_2(\F_2)\) & vectorization & Linear bijection after cell-order choice; basis/order dependent.\\
\addlinespace
Four coefficients & \(c\in\F_2^4\) & cyclic-basis identification & Vector-space bijection to \(A\); does not inherit ordinary \(2\times2\) matrix multiplication.\\
\addlinespace
Chain-ring scalar & \(a\in A\cong\F_2[u]/(u^4)\) & regular representation \(\Lambda\) & Injective algebra representation into \(M_4(\F_2)\) after multiplication on \(A\) is chosen.\\
\addlinespace
One-system state & \(\psi\in A^n\) & effect application & Free \(A\)-module object; one CM is only one scalar of \(A\), not a whole state.\\
\addlinespace
Bipartite resource & \(M\in M_{m,n}(A)\) & \(M\mapsto UMV^T\) & Local \(A\)-linear action; distinct type from original Boolean CMs.\\
\addlinespace
Effect & \(e\in (A^n)^*\cong A^{1\times n}\) & \(e(\psi)=e\psi\) & \(A\)-linear functional; row identification uses a free-module basis.\\
\addlinespace
Measurement & rows of \(B\in\GL_n(A)\) & \(B\psi\) & Reversible coordinate change; additional modal choice.\\
\addlinespace
Outcome scalar & \(e\psi\in A\) or \(eMf^T\in A\) & \(a\mapsto[a\neq0]\) & Support zero-test; not a ring/semiring homomorphism in general.\\
\addlinespace
Symbolic enriched layer & \(\Sring=\B\otimes_{\F_2}A\) & \(v_A=v\otimes\Id_A\) & Canonical once \(A\) is chosen; preserves ring, module, matrix, and tensor operations.\\
\addlinespace
Symbolic modal outcome & \(q\in\Sring\) & \(\Supp(q)\in\B\) & Exact valuation-wise zero/nonzero semantics; support itself is not an algebra homomorphism.\\
\bottomrule
\end{longtable}

A key typing correction is
\[
\boxed{\text{one }2\times2\text{ Boolean CM}\neq\text{one modal state}.}
\]
After a four-feature basis is chosen, one CM becomes one scalar of \(A\). A one-bit modal state needs
two such scalars:
\[
\psi=a_0\lvert0\rangle+a_1\lvert1\rangle\in A^2.
\]
A two-party resource requires four such scalars.

\section{Q1: What is a formula-valued LM mathematically?}

Let \(\B\) be the Lindenbaum Boolean algebra of propositional formulas modulo logical equivalence.
Write it as a Boolean ring using
\[
f+g=f\oplus g,\qquad fg=f\land g,\qquad \neg f=1+f.
\]
For freely assignable symbols \(X,Y\), define
\[
X_1=X,\quad X_2=\neg X,\qquad
Y_1=Y,\quad Y_2=\neg Y.
\]
For a binary truth table \(C=(c_{ij})\in M_2(\F_2)\), Paper B defines
\[
\LM_{X,Y}(C)
=
\bigoplus_{i,j}c_{ij}\lvert X_i\rangle\langle Y_j\rvert
\in M_2(\B).
\]
Equivalently, after expansion,
\[
\LM_{X,Y}(\Theta)
=
\begin{pmatrix}
X\,\Theta\,Y & X\,\Theta\,\neg Y\\
\neg X\,\Theta\,Y & \neg X\,\Theta\,\neg Y
\end{pmatrix}.
\]

\begin{theorem}[Algebraic status of the Paper-B LM family]
For fixed freely assignable \(X,Y\), the map
\[
\ell_{X,Y}:M_2(\F_2)_{\mathrm{coeff}}\longrightarrow M_2(\B),
\qquad C\longmapsto\LM_{X,Y}(C),
\]
is injective and \(\F_2\)-linear. It preserves entrywise Boolean operations. It is not, in general,
a homomorphism from ordinary Boolean-ring matrix multiplication on \(M_2(\F_2)\) to matrix
multiplication in \(M_2(\B)\).
\end{theorem}

\begin{proof}
Linearity follows from the XOR sum defining the LM. Positive valuation sends the LM back to its
coefficient CM, so equality of two LMs implies equality of their coefficient matrices. Therefore the
map is injective. Entrywise Boolean operations commute with the substitution producing the four
polarity positions. Ordinary matrix product is different: multiplication mixes row/column positions
by XOR--AND contraction. Exhaustive evaluation of all \(16^2=256\) products of the 16 Paper-B
LMs shows that 144 products leave the restricted 16-element LM family, so the family is not closed
under this multiplication.
\end{proof}

\begin{remark}
The most economical description is therefore: a Paper-B LM is a special predicate-valued matrix
over a free/Lindenbaum Boolean ring. The ambient space \(M_2(\B)\) is natural; the 16 matrices
generated by binary truth functions are a distinguished subset/subalgebra only for the entrywise
Boolean operations. Calling the construction merely a relation would lose the XOR--AND scalar
algebra used by Paper B.
\end{remark}

\section{Q2: What is LM pairing mathematically?}

For a formula \(F\in\B\), Paper B defines
\[
s(F)=
\begin{pmatrix}F\\ \neg F\end{pmatrix},
\qquad
s(F)^T=(F,\neg F).
\]
The symbolic pairing is
\[
P_C(L_1,L_2)
=
s(L_1)^T\LM_{X,Y}(C)s(L_2),
\]
where scalar multiplication is conjunction and scalar addition is XOR.

\begin{proposition}[LM pairing is Boolean-ring matrix contraction]
For the special LM associated with a binary connective \(\Theta\),
\[
\langle L_1\vert\LM_{X,Y}(\Theta)\vert L_2\rangle
\equiv
(L_1\leftrightarrow X)\;\Theta\;(Y\leftrightarrow L_2).
\]
Thus the Paper-B ``logical pairing'' is ordinary matrix contraction over the Boolean ring \(\B\),
with a non-linear special state embedding \(F\mapsto(F,\neg F)\).
\end{proposition}

The roles of the connectives are exact:
\begin{itemize}
\item AND is scalar multiplication;
\item XOR is scalar addition and reduction;
\item negation appears in the complement-paired state encoding;
\item OR is not the matrix-addition operation used by the pairing.
\end{itemize}

The pairing may be described semantically as predicate evaluation or relationship extraction, but
those descriptions do not add algebraic structure beyond the Boolean-ring contraction.

\section{Q3: Valuation naturality}

Let
\[
v:\B\longrightarrow\F_2
\]
be a Boolean valuation. Extend \(v\) entrywise to matrices and vectors.

\begin{theorem}[Valuation naturality]
For compatible Boolean-ring matrices and vectors,
\[
v(MN)=v(M)v(N),
\]
and
\[
v(x^TMy)=v(x)^T\,v(M)\,v(y).
\]
The same statement holds for finite Kronecker tensor products.
\end{theorem}

\begin{proof}
Every output entry is a finite expression built from Boolean-ring addition and multiplication. A
Boolean valuation is a homomorphism for those operations, so it may be moved through each finite
sum and product.
\end{proof}

This theorem is already present in Paper B for the LM pairing case. Its broader matrix/tensor version
is a routine homomorphism-extension fact. It is useful for organizing the bridge, but it should not be
presented as deep novelty.

\section{A decisive negative theorem: direct valuation cannot produce chain-ring amplitudes}

The strongest possible bridge would send Boolean formulas directly into the later amplitude ring
\[
A=\F_2[u]/(u^4).
\]
That is impossible except for the embedded Boolean scalars.

\begin{theorem}[Idempotent obstruction]
Let \(\B\) be a Boolean ring and let \(h:\B\to A\) be a unital ring homomorphism. Then every
\(h(b)\) is idempotent. In the local chain ring \(A=\F_2[u]/(u^4)\), the only idempotents are
\(0\) and \(1\). Therefore
\[
\operatorname{im}h\subseteq\{0,1\}\cong\F_2,
\]
and \(h\) factors through an ordinary Boolean valuation
\[
\B\longrightarrow\F_2\hookrightarrow A.
\]
\end{theorem}

\begin{proof}
Every element of a Boolean ring satisfies \(b^2=b\). Hence
\[
h(b)^2=h(b^2)=h(b).
\]
A local ring has no idempotents other than \(0\) and \(1\). The stated chain ring is local, so the
image is Boolean.
\end{proof}

\begin{corollary}
Nilpotent and nontrivial unit amplitudes such as
\[
u,\quad u^2,\quad u^3,\quad 1+u
\]
can never arise as direct values of Paper-B Boolean formulas under a ring-homomorphic valuation.
\end{corollary}

This rules out the claim that the chain-ring amplitude theory is obtained by simply changing the
codomain of Paper-B valuation.

\section{The CM-to-chain-ring identification is linear before it is multiplicative}

There is a natural-looking sequence
\[
M_2(\F_2)
\xrightarrow{\vecop_{\kappa}}
\F_2^4
\xrightarrow{\alpha_{\beta}}
A,
\]
where \(\kappa\) fixes a cell order and \(\beta\) fixes an ordered \(\F_2\)-basis of \(A\).
Both arrows are vector-space bijections.

They are not an identification of ordinary \(2\times2\) matrix multiplication with the chain-ring
multiplication. Two immediate invariants separate the algebras:
\begin{itemize}
\item \(M_2(\F_2)\) is noncommutative, while \(A\) is commutative;
\item \(|\GL_2(\F_2)|=6\), while \(A\) has 8 units.
\end{itemize}

Therefore the later cyclic product
\[
E_r\star E_s=E_{r+s\bmod4}
\]
is additional algebraic structure placed on the four-component coefficient space. Once supplied,
\[
A\cong\F_2[C_4]\cong\F_2[u]/(u^4),
\]
and its regular representation
\[
\Lambda:A\hookrightarrow M_4(\F_2)
\]
is an injective algebra homomorphism. The multiplication itself, however, is not forced by the
original LM valuation theorem.

\section{The correct positive bridge: scalar enrichment}

The direct-valuation obstruction suggests the right construction.

\begin{definition}[Enriched symbolic scalar algebra]
Let
\[
\Sring=\B\otimes_{\F_2}A.
\]
For the current chain ring,
\[
\Sring\cong \B[u]/(u^4).
\]
\end{definition}

\begin{theorem}[Canonical specialization after choosing \(A\)]
For every Boolean valuation \(v:\B\to\F_2\), there is a canonical \(A\)-algebra specialization
\[
v_A=v\otimes\Id_A:\Sring\to A,
\]
defined on pure tensors by
\[
v_A(f\otimes a)=v(f)a.
\]
It preserves addition and multiplication and therefore extends entrywise to matrices, vectors,
contractions, and tensor products. In particular,
\[
v_A(e\psi)=v_A(e)v_A(\psi),
\]
and
\[
v_A(E M F^T)=v_A(E)v_A(M)v_A(F)^T.
\]
\end{theorem}

\begin{remark}[Universal property]
The tensor product is the coproduct of the commutative \(\F_2\)-algebras \(\B\) and \(A\).
Therefore this is a genuine universal property, but it is a universal property \emph{after} the amplitude
algebra \(A\) has been independently selected. It does not derive \(A\) from the formula calculus.
\end{remark}

For finitely many Boolean generators, if
\[
\Omega=\operatorname{Hom}_{\mathrm{Bool}}(\B,\Bool)
\]
is the valuation set, then extensionally
\[
\B\cong\F_2^{\Omega}
\quad\text{and hence}\quad
\boxed{\Sring\cong A^{\Omega}.}
\]
So an enriched symbolic scalar may be read simply as an \(A\)-valued function of the Boolean
valuation. This is the finite Boolean-power viewpoint.

\section{The strongest bridge theorem: symbolic support realizes modal possibility}

Let
\[
q\in\Sring\cong A^{\Omega}.
\]
Define the symbolic support predicate \(\Supp(q)\in\B\) to have exactly the valuation set
\[
\{v\in\Omega: v_A(q)\neq0\}.
\]
For
\[
q=f_0+f_1u+f_2u^2+f_3u^3,
\qquad f_i\in\B,
\]
one may write
\[
\Supp(q)=f_0\lor f_1\lor f_2\lor f_3.
\]
Although this coordinate formula uses the displayed basis, the resulting Boolean predicate is
basis-independent up to logical equivalence because it is defined by the intrinsic condition
\(v_A(q)\neq0\).

\begin{theorem}[Symbolic-support bridge]
Let
\[
e\in\Sring^{1\times n},\qquad \psi\in\Sring^n.
\]
Define
\[
\Poss(e,\psi)=\Supp(e\psi)\in\B.
\]
Then for every Boolean valuation \(v\),
\[
\boxed{
v\bigl(\Poss(e,\psi)\bigr)=1
\iff
v_A(e)\,v_A(\psi)\neq0.}
\]
For a bipartite symbolic resource \(M\),
\[
\boxed{
v\left(\Supp(eMf^T)\right)=1
\iff
v_A(e)v_A(M)v_A(f)^T\neq0.}
\]
\end{theorem}

\begin{proof}
Choose an \(\F_2\)-basis \(\{\beta_k\}\) of \(A\) and write the symbolic scalar contraction as
\[
q=\sum_k f_k\beta_k.
\]
At valuation \(v\),
\[
v_A(q)=\sum_k v(f_k)\beta_k.
\]
Linear independence of the basis gives
\[
v_A(q)=0\iff v(f_k)=0\ \text{for all }k.
\]
Therefore
\[
v_A(q)\neq0
\iff
\bigvee_kv(f_k)=1
\iff
v\left(\bigvee_k f_k\right)=1.
\]
\end{proof}

\begin{corollary}[Satisfiability/possibility bridge]
\[
\Supp(e\psi)\text{ is satisfiable}
\iff
\exists v:\;v_A(e)v_A(\psi)\neq0.
\]
The same holds for bipartite contractions \(eMf^T\).
\end{corollary}

This is the strongest bridge found in the investigation. It says that a single symbolic contraction can
compute an ordinary Boolean formula describing exactly the parameter valuations under which the
corresponding concrete ring-modal event is possible.

\section{Q4: Can formula-valued LMs represent modal effects?}

Two meanings of ``represent'' must be distinguished.

\subsection{As measurement semantics: only a restricted subset}

The Paper-B logical state map is
\[
s(L)=\begin{pmatrix}L\\\neg L\end{pmatrix}.
\]
For every valuation,
\[
v(s(L))\in
\left\{
\begin{pmatrix}1\\0\end{pmatrix},
\begin{pmatrix}0\\1\end{pmatrix}
\right\}.
\]
So the bare Paper-B bra/ket construction yields only coordinate rows after valuation. It does not
directly yield the embedded-\(\F_2\) effect \((1,1)\), which already occurs in a standard modal
\(X\)-type basis.

\begin{proposition}[Negative measurement-family characterization]
The original Paper-B LM pairing semantics does not generate the complete 192 reversible projective
modal bases over \(A\). It does not even generate every reversible basis in the embedded
\(\F_2\) subtheory.
\end{proposition}

This is a negative version of the requested measurement-family characterization.

\subsection{As coefficient data: every effect can be encoded after choices}

Fix a CM cell ordering and the cyclic \(\F_2\)-basis of \(A\). Every \(a\in A\) has a unique four-bit
coefficient representation, hence a unique Boolean CM in that basis, and therefore a Paper-B LM
representative in a chosen operand frame. An effect
\[
e=(a,b)\in A^2
\]
can consequently be encoded by two such LMs, and a \(2\times2\) reversible basis by four.

Thus every modal effect is LM-encodable as coefficient data, but this encoding is:
\begin{itemize}
\item dependent on the chosen four-feature basis;
\item dependent on CM cell order;
\item dependent on the chosen LM operand frame/positive-valuation convention;
\item not a derivation of effect semantics;
\item not an identification of LM matrix multiplication with \(A\)-multiplication.
\end{itemize}

The correct statement is therefore:
\begin{quote}
Every modal effect is LM-encodable after the chain-ring coefficient basis is fixed, but not every modal
effect is an LM effect in the original Paper-B pairing sense.
\end{quote}

\section{Q5: ``nonzero = possible'' is extra structure}

The later modal source explicitly introduces nonzero readout as an additional contract. The algebra
shows why it cannot be hidden inside ordinary Boolean valuation semantics.

Define
\[
\sigma:A\to\{0,1\},\qquad
\sigma(a)=\begin{cases}0,&a=0,\\1,&a\neq0.\end{cases}
\]
This map is neither additive in the \(\F_2\) sense nor multiplicative to Boolean support.
For example,
\[
\sigma(u)=\sigma(u^2)=1,
\]
but
\[
\sigma(u+u^2)=1\neq1\oplus1=0.
\]
More importantly, zero divisors give
\[
u\neq0,\qquad u^3\neq0,\qquad uu^3=u^4=0,
\]
so
\[
\sigma(u)\land\sigma(u^3)=1
\quad\text{while}\quad
\sigma(uu^3)=0.
\]

\begin{proposition}[Support is not a semiring homomorphism]
The nonzero map \(\sigma:A\to\Bool\) cannot be a homomorphism from the chain ring to either the
Boolean ring \((\Bool,\oplus,\land)\) or the ordinary Boolean semiring \((\Bool,\lor,\land)\).
Consequently, the support semantics is an independently selected zero-test rather than a functorial
quotient of the amplitude algebra.
\end{proposition}

The symbolic-support theorem does not contradict this. It says that, \emph{once the zero-test is chosen},
its valuation fibers can be represented exactly by Boolean formulas.

\section{Q6: Symbolic observables/effects}

The most useful symbolic generalization is not a new Boolean algebra of observables but a module of
symbolic effects:
\[
\mathcal E_n=(\Sring^n)^*\cong\Sring^{1\times n}.
\]
An effect may carry genuine formulas in its coefficients, for example
\[
e=
\bigl(
(x\land\neg y)+u(x\oplus z),\;
 u^2(f\leftrightarrow g)+u^3z
\bigr).
\]
For a symbolic state \(\psi\), the contraction
\[
e\psi\in\Sring
\]
retains the complete symbolic \(A\)-valued amplitude. Applying support gives
\[
\Supp(e\psi)\in\B,
\]
a Boolean formula describing exactly when the effect is modally possible after specialization.

Thus the framework naturally exposes two levels of output:
\[
\boxed{e\psi\ \text{(symbolic amplitude)}}
\qquad\text{and}\qquad
\boxed{\Supp(e\psi)\ \text{(symbolic possibility condition)}}.
\]

Calling this an ``effect algebra'' would currently be premature because no partial sum, orthosupplement,
or effect-algebra order has been shown to satisfy the standard axioms. ``Symbolic effect module'' or
``predicate-parametrized effect calculus'' is more accurate.

\section{Q7: Sequential measurement and update}

The chain ring creates a normalization problem because possible amplitudes can be nonunits. However,
reversible-basis measurements admit a clean branch update that requires no division by the observed
amplitude.

Let
\[
B\in\GL_n(A)
\]
be a reversible measurement basis, with rows interpreted as effects. Let \(P_i\) be the diagonal
coordinate projector having a single \(1\) in coordinate \(i\). Define
\[
\boxed{J_i^B=B^{-1}P_iB.}
\]

\begin{theorem}[Division-free reversible-basis update]
For every commutative ring \(A\) and every \(B\in\GL_n(A)\), the maps
\(J_i^B\) satisfy
\[
(J_i^B)^2=J_i^B,
\]
\[
J_i^BJ_j^B=0\quad(i\neq j),
\]
and
\[
\sum_iJ_i^B=I.
\]
If
\[
c_i=(B\psi)_i,
\]
then
\[
\boxed{J_i^B\psi\neq0\iff c_i\neq0.}
\]
Moreover,
\[
B(J_i^B\psi)=P_iB\psi,
\]
so a repeat of the same measurement has support only on outcome \(i\).
\end{theorem}

\begin{proof}
Idempotence, orthogonality, and resolution of the identity follow by conjugating the coordinate
projectors. Let
\[
b_i=B^{-1}e_i.
\]
Then
\[
J_i^B\psi=c_i b_i.
\]
The column \(b_i\) is unimodular because it is a column of an invertible matrix. Hence its entries
generate the unit ideal, so there exist \(r_j\) with
\[
\sum_j r_j(b_i)_j=1.
\]
If \(c_i b_i=0\), multiplying by the \(r_j\) combination yields \(c_i=0\). The converse is
immediate. Finally,
\[
B J_i^B\psi=P_iB\psi.
\]
\end{proof}

\subsection{Projective well-definedness}

If the effect rows are independently rescaled by units,
\[
B'=DB
\]
with \(D\) diagonal and invertible, then
\[
(B')^{-1}P_iB'
=B^{-1}D^{-1}P_iDB
=B^{-1}P_iB.
\]
Therefore the branch projector is invariant under the same row-unit projectivization already used by
the modal support theory. Outcome permutations only relabel \(i\).

\subsection{Symbolic update}

For
\[
\widetilde B\in\GL_n(\Sring)
\]
define
\[
\widetilde J_i=\widetilde B^{-1}P_i\widetilde B.
\]
Then every specialization satisfies
\[
v_A(\widetilde J_i)=J_i^{v_A(\widetilde B)}.
\]
So valuation commutes with the branch-update construction as long as the symbolic basis remains
invertible under the relevant specialization.

\section{Why this is not yet a closed process theory}

The projector update is algebraically clean, but the zero divisors prevent the set of all nonzero
states from being a well-behaved monoidal state class.

\subsection{Nonzero states are not tensor-closed}

Take
\[
\psi=(u,0),\qquad \phi=(u^3,0).
\]
Both are nonzero, but over \(A\)
\[
\psi\otimes_A\phi=0
\]
because \(uu^3=u^4=0\).

\subsection{Unimodular preparation is not conditioning-closed}

Restricting preparations to unimodular vectors repairs the elementary tensor-annihilation problem for
simple tensors, but conditioning can leave that class. Consider the bipartite coefficient matrix
\[
M=\begin{pmatrix}1&0\\0&u\end{pmatrix}.
\]
As a four-component vector it is unimodular because one entry is a unit. Applying Alice's second
standard effect leaves Bob in
\[
(0,u),
\]
which is nonzero but nonunimodular.

Therefore a satisfactory closed sequential theory still needs an explicit decision about admissible
states, branch objects, quotients, ideals, or a richer process semantics. The symbolic predicate layer
can describe when a branch exists, but it does not turn a nonunit branch amplitude into an invertible
one.

\section{Failure results and exact boundaries}

\begin{longtable}{p{0.34\textwidth}p{0.58\textwidth}}
\toprule
\textbf{Proposed identification} & \textbf{Result} \\
\midrule
\endhead
``An LM is a modal effect.'' & False as stated. Paper-B bras evaluate only to coordinate effects.\\
\addlinespace
``Boolean valuation directly gives \(A\)-amplitudes.'' & False. Every Boolean-ring homomorphism into the local chain ring lands in \(\{0,1\}\).\\
\addlinespace
``The ordinary \(2\times2\) CM algebra is the phase ring.'' & False. The four-component identification is linear; the cyclic product is additional structure.\\
\addlinespace
``Valuation naturality derives modal possibility.'' & False. It derives concrete amplitudes. The zero/nonzero interpretation is a separate support contract.\\
\addlinespace
``Nonzero support is ordinary relational semantics.'' & False compositionally. XOR cancellation and zero divisors prevent support from being a semiring homomorphism.\\
\addlinespace
``Existing LM semantics generates all 192 measurements.'' & False. The full family is LM-encodable as coefficient data after adding the \(A\)-structure, but not generated by the original complement-paired effect semantics.\\
\addlinespace
``Symbolic formulas solve normalization.'' & False. They can describe branch conditions but do not invert nonunits.\\
\addlinespace
``All nonzero ring states form a monoidal state theory.'' & False: \(u\otimes_Au^3=0\).\\
\addlinespace
``Unimodular preparations solve closure.'' & False: conditioning can yield nonunimodular nonzero states.\\
\addlinespace
``The bridge extends unchanged to arbitrary finite rings.'' & False as stated. The clean \(\B\otimes_{\F_2}A\) enrichment requires an \(\F_2\)-algebra such as a characteristic-two ring; other scalar systems require a modified control/Boolean-power construction.\\
\bottomrule
\end{longtable}

\section{Prior-art matrix}

The novelty search translated each candidate theorem into conventional mathematics rather than
searching only CM terminology.

\begin{longtable}{p{0.31\textwidth}p{0.62\textwidth}}
\toprule
\textbf{Construction/result} & \textbf{Prior-art assessment} \\
\midrule
\endhead
\(2\times2\) matrices for the 16 binary truth functions & Known. Bricken's 1997 notes explicitly discuss XOR as binary addition, AND as multiplication, all 16 binary logic matrices, truth-value vectors, and expressions of the form \(\langle x\vert B\vert y\rangle\). Edwards, Stern, and matrix-logic traditions are older.\\
\addlinespace
Truth vectors and matrix logical gates & Known. Mizraji's vector logic represents classical logical connectives by matrices acting on Kronecker products of truth vectors.\\
\addlinespace
Quantum-flavored logical observables & Known adjacent work. Eigenlogic treats logical observables in a genuine vector/Hilbert-space style distinct from Paper-B formula-valued LMs.\\
\addlinespace
Formula/Lindenbaum algebra modulo logical equivalence & Standard algebraic logic.\\
\addlinespace
Valuation commuting with matrix expressions & Standard homomorphism preservation and already proved in Paper B for LM pairing.\\
\addlinespace
\(\B\otimes A\) enrichment and \(A^\Omega\) interpretation & Standard base-change / Boolean-power mathematics. The contribution is its application as the exact typed bridge between the two CM layers.\\
\addlinespace
Boolean-to-local-ring idempotent obstruction & Elementary ring theory; not a novelty target.\\
\addlinespace
Modal effects with possible iff nonzero & Established in modal quantum theory over finite fields, especially Schumacher--Westmoreland.\\
\addlinespace
Finite-field/XOR quantum-like programming & Established adjacent work in James--Ortiz--Sabry.\\
\addlinespace
Semiring-valued quantum-like process theories & Established broad area; Gogioso's framework is a direct comparison point.\\
\addlinespace
Possibilistic contextuality as support/global-section obstruction & Established by Abramsky--Brandenburger and follow-on sheaf-theoretic work.\\
\addlinespace
States/effects and categorical measurement/update & Established broader categorical foundations, including effectus theory.\\
\addlinespace
Exact symbolic-support theorem for the CM-derived chain-ring enrichment & No exact prior formulation located in this search. Best characterized as a potentially useful new synthesis theorem rather than a new base algebraic construction.\\
\addlinespace
Negative theorem that original LM pairing does not generate the 192-basis family & Appears to be a useful internal comparison theorem; mathematically elementary once the types are separated.\\
\addlinespace
\(B^{-1}P_iB\) update & Standard conjugated-projector/change-of-basis algebra; useful as the natural division-free update for the current modal model.\\
\bottomrule
\end{longtable}

\section{Adversarial interpretation}

The project survives an adversarial reading, but in a narrower form than a claim of derivation.

\subsection{What is routine}
\begin{itemize}
\item valuation naturality is homomorphism preservation;
\item matrix semantics for logic has substantial prior art;
\item semiring/ring-valued process models and modal nonzero semantics have prior art;
\item conjugated coordinate projectors are standard algebra.
\end{itemize}

\subsection{What is genuinely informative in the synthesis}
\begin{itemize}
\item the idempotent obstruction proves that the strongest direct-valuation bridge is impossible;
\item the tensor-product scalar enrichment gives the correct universal replacement;
\item the support predicate gives an exact formula-level representation of concrete modal possibility;
\item the effect-family comparison identifies which sense of ``LM measurement'' does and does not extend to the 192 modal bases;
\item the update construction shows that nonunit amplitudes do not by themselves block a division-free branch map;
\item the zero-divisor closure counterexamples locate the remaining process-theoretic weakness precisely.
\end{itemize}

\section{Computational verification}

A fresh independent checker is included as
\begin{center}
\texttt{tests/lm\_modal\_bridge\_checks.py}.
\end{center}
It uses only the Python standard library. The exact machine-readable output is included as
\begin{center}
\texttt{tests/lm\_modal\_bridge\_checks\_output.json}.
\end{center}

The expanded run verified:

\begin{longtable}{p{0.59\textwidth}r}
\toprule
\textbf{Check} & \textbf{Count/result} \\
\midrule
\endhead
Elements of \(A=\F_2[u]/(u^4)\) & 16\\
Units & 8\\
Idempotents & exactly \(0,1\)\\
\(|\GL_2(A)|\) & 24,576\\
Unordered reversible projective two-outcome bases & 192\\
Projector idempotence checks & 384\\
Projector orthogonality checks & 384\\
Resolution-of-identity checks & 192\\
Branch nonzero iff effect coefficient nonzero & 98,304\\
Same-basis repeatability checks & 98,304\\
Unit row-scaling invariance checks & 24,576\\
Symbolic support fiber checks & 262,144\\
Paper-B LM matrix products tested & 256\\
Products remaining in restricted 16-LM family & 112\\
Products leaving restricted 16-LM family & 144\\
Nonzero-state tensor annihilation example & confirmed\\
Unimodular preparation to nonunimodular conditional example & confirmed\\
\bottomrule
\end{longtable}

The 262,144 symbolic-support checks exhaust all
\[
16^4=65,536
\]
symbolic \(A\)-valued scalars whose four coefficient formulas are Boolean functions of two variables,
and all four valuations of those variables.

The package also preserves the original checker and output generated during the first pass, plus
rerun logs showing that both the original and expanded scripts execute successfully in the current
environment.

\section{Paper implications}

The bridge is strong enough to deserve a substantial theory section, but the present evidence does
not support a standalone paper whose headline claim is that Paper-B LMs derive modal quantum
measurement.

The best placement is a central section of the broader Boolean--modal synthesis paper. A strong
narrative is:
\[
\text{formula LM}
\to
\text{Boolean-ring contraction}
\to
\text{valuation naturality}
\to
\boxed{\text{idempotent obstruction}}
\to
\text{\(A\)-enrichment}
\to
\boxed{\text{symbolic-support theorem}}
\to
\text{modal effects/bases}
\to
\boxed{\text{extra support contract}}
\to
\text{division-free branch update}
\to
\boxed{\text{tensor/update closure obstruction}}.
\]

This structure has three advantages:
\begin{enumerate}
\item it answers the skeptical question of why the original logical pairing has anything to do with the
later modal measurement rule;
\item it clearly separates generic prior art from CM-specific typing and synthesis;
\item it uses negative theorems as part of the contribution rather than forcing an overstated positive
unification.
\end{enumerate}

\section{Strongest defensible theorem and final answer}

The single best theorem to carry forward is the support-bridge statement
\[
\boxed{
v\!\left(
\Supp\bigl(\langle E,S\rangle_{\B\otimes A}\bigr)
\right)=1
\iff
\left\langle
v_A(E),v_A(S)
\right\rangle_A\neq0.}
\]
Together with its existential form,
\[
\boxed{
\Supp(\langle E,S\rangle)
\text{ satisfiable}
\iff
\exists v:\langle v_A(E),v_A(S)\rangle_A\neq0.}
\]

It should be stated beside the negative companion theorem
\[
\boxed{
\text{No direct Boolean-ring valuation }
\B\to\F_2[u]/(u^4)
\text{ can generate non-Boolean amplitudes.}}
\]

The pair identifies the bridge almost exactly:
\begin{quote}
The formula-valued LM calculus and the later modal theory do more than merely coexist, but the modal
theory is an enriched realization of the symbolic Boolean semantics, not a theorem forced by the
original LM semantics alone. The correct bridge is scalar enrichment plus specialization, followed by
an independently declared support interpretation and measurement-family choice.
\end{quote}

\appendix
\section{Reproducibility notes}

Run the expanded checker from the root of this package with
\begin{lstlisting}[language=bash]
python tests/lm_modal_bridge_checks.py
\end{lstlisting}
The output should match \texttt{tests/lm\_modal\_bridge\_checks\_output.json} modulo the environment
string. The script contains no network calls and depends only on the Python standard library.

The prior adversarial-review archive is preserved under \texttt{inputs/} and unpacked under
\texttt{prior\_adversarial\_review/}. The package manifest records SHA-256 hashes for every included
file.

\section{Search strategy}

The prior-art search used conventional mathematical translations rather than CM-specific vocabulary.
Queries included Boolean matrix logic, matrix logic and vector logic, formula/Lindenbaum algebra,
matrices over semirings, Boolean powers, modal quantum theory, finite-field quantum computation,
possibilistic contextuality, effectus theory, categorical quantum mechanics, and semiring-valued
quantum-like theories. The machine-readable search ledger is included separately in
\texttt{research/SEARCH\_LOG.csv}.

\begin{thebibliography}{99}

\bibitem{Bricken1997}
William Bricken.
\newblock \emph{Notes on Matrix Techniques for Logic}.
\newblock March 1997.
\newblock \url{https://wbricken.com/pdfs/01bm/01math/03math-supporting/math-tangential/04matrix-tech.pdf}.

\bibitem{Edwards1972}
C. R. Edwards.
\newblock The logic of Boolean matrices.
\newblock \emph{The Computer Journal}, 15(3):247--253, 1972.
\newblock doi:10.1093/comjnl/15.3.247.

\bibitem{Stern1988}
August Stern.
\newblock \emph{Matrix Logic: Theory and Applications}.
\newblock North-Holland, 1988.
\newblock doi:10.1016/C2009-0-14445-8.

\bibitem{Mizraji1992}
Eduardo Mizraji.
\newblock Vector logics: The matrix-vector representation of logical calculus.
\newblock \emph{Fuzzy Sets and Systems}, 50(2):179--185, 1992.
\newblock doi:10.1016/0165-0114(92)90216-Q.

\bibitem{Mizraji2008}
Eduardo Mizraji.
\newblock Vector logic: A natural algebraic representation of the fundamental logical gates.
\newblock \emph{Journal of Logic and Computation}, 18(1):97--121, 2008.
\newblock doi:10.1093/logcom/exm057.

\bibitem{ChengQi2010}
Daizhan Cheng and Hongsheng Qi.
\newblock A linear representation of dynamics of Boolean networks.
\newblock \emph{IEEE Transactions on Automatic Control}, 55(10):2251--2258, 2010.
\newblock doi:10.1109/TAC.2010.2043294.

\bibitem{Cheng2011}
Daizhan Cheng, Yin Zhao, and Xiangru Xu.
\newblock Matrix approach to Boolean calculus.
\newblock In \emph{Proceedings of the 50th IEEE CDC-ECC}, pages 6950--6955, 2011.
\newblock doi:10.1109/CDC.2011.6160289.

\bibitem{DuboisToffano2016}
Francois Dubois and Zeno Toffano.
\newblock Eigenlogic: A quantum view for multiple-valued and fuzzy systems.
\newblock arXiv:1607.03509, 2016.

\bibitem{SchumacherWestmoreland2012}
Benjamin Schumacher and Michael D. Westmoreland.
\newblock Modal quantum theory.
\newblock \emph{Foundations of Physics}, 42:918--925, 2012.
\newblock doi:10.1007/s10701-012-9650-z; arXiv:1010.2929.

\bibitem{JamesOrtizSabry2011}
Roshan P. James, Gerardo Ortiz, and Amr Sabry.
\newblock Quantum computing over finite fields: Reversible relational programming with exclusive disjunctions.
\newblock arXiv:1101.3764, 2011.

\bibitem{AbramskyBrandenburger2011}
Samson Abramsky and Adam Brandenburger.
\newblock The sheaf-theoretic structure of non-locality and contextuality.
\newblock \emph{New Journal of Physics}, 13(11):113036, 2011.
\newblock doi:10.1088/1367-2630/13/11/113036; arXiv:1102.0264.

\bibitem{ChoEtAl2015}
Kenta Cho, Bart Jacobs, Bas Westerbaan, and Abraham Westerbaan.
\newblock An introduction to effectus theory.
\newblock arXiv:1512.05813, 2015.

\bibitem{Gogioso2017}
Stefano Gogioso.
\newblock Fantastic quantum theories and where to find them.
\newblock arXiv:1703.10576, 2017.

\bibitem{PaperB}
Brian Theory.
\newblock \emph{Operator-Level Boolean Computation with Correspondence Matrices}.
\newblock Unpublished manuscript supplied in the ChatGPT File Library, accessed 18 September 2026.

\bibitem{ModalAddendum}
Brian Theory.
\newblock \emph{Pure Boolean Modal CM Addendum}.
\newblock Research manuscript/source file \texttt{PURE\_BOOLEAN\_MODAL\_CM\_ADDENDUM.tex}, accessed 18 September 2026.

\end{thebibliography}

\end{document}
